\chapter{Antecedentes}
\label{ch:antecedentes}

\paragraph*{}El análisis de supervivencia es una técnica estadística diseñada para estudiar el tiempo que transcurre hasta que ocurre un evento determinado. Su característica más valiosa es que permite trabajar con datos censurados, es decir, con casos en los que el evento todavía no se ha producido durante el periodo de observación, algo que otros métodos suelen ignorar. En medicina se usa para estudiar la evolución de pacientes y en ingeniería para estimar la vida útil de componentes. En educación, aplicado al abandono universitario, resulta especialmente útil porque no solo permite saber cuántos estudiantes abandonan, sino también cuándo lo hacen y bajo qué circunstancias es más probable que ocurra e

\paragraph*{}Dentro de este campo, los modelos más consolidados son Kaplan-Meier, que estima la probabilidad de permanencia a lo largo del tiempo, el Test de Log-Rank, que permite comparar trayectorias entre grupos, y la regresión de Cox, que identifica qué variables ---rendimiento académico, edad, situación socioeconómica, entre otras--- influyen de forma significativa en el riesgo de abandono. Más recientemente, han ganado terreno enfoques como los modelos paramétricos de Weibull o los Random Survival Forests, que incorporan la potencia del aprendizaje automático para mejorar la capacidad predictiva del análisis.

\paragraph*{}La literatura sobre abandono universitario aborda el problema desde distintos ángulos. Algunos trabajos se centran en los factores personales, académicos o institucionales que condicionan la trayectoria del estudiante. Otros apuntan directamente a la detección temprana de perfiles de riesgo. Esta segunda línea es la que más se aproxima al enfoque de este proyecto, donde el análisis no es un fin en sí mismo sino una herramienta para anticipar y prevenir.

\paragraph*{}El antecedente más directo es el Trabajo Fin de Grado de Lucía Cabezuelo Pérez (2025), titulado \textit{``Análisis, Diseño e Implementación de un Dashboard Interactivo para el Análisis de Supervivencia aplicado al Abandono Escolar Universitario''}. En él se desarrolló una aplicación web en Python y Dash que integraba los modelos de Kaplan-Meier, regresión de Cox y Test de Log-Rank, con visualizaciones interactivas y un módulo de inteligencia artificial basado en Llama 3 a través de Ollama, capaz de generar interpretaciones automáticas de los resultados. Fue un trabajo sólido, que demostró la viabilidad del enfoque y dejó una base técnica sobre la que merece la pena seguir construyendo.

\paragraph*{}Este proyecto recoge ese testigo y lo lleva más lejos. Se amplía el abanico de modelos estadísticos disponibles, se incorporan técnicas de aprendizaje automático aplicadas a la supervivencia, se sustituye el módulo de IA por una solución más rápida y eficiente, y se enriquece la plataforma con soporte multilingüe, exportación de resultados y una interfaz rediseñada. No se trata de rehacer lo que ya funcionaba, sino de completar lo que faltaba.
